\chapter{Parallel programming}
\label{section:parallel-programming}

Peano has originally been written for MPI+X on CPU-based architectures.
On the MPI side, I provide some wrappers around mainstream MPI routines which
augment the MPI things with deadlock detection, explicit buffering of unexpected
messages or global synchronisation.
On the X-side, \Peano\ has its own 


\section{Shared memory}

\subsection{Protecting shared memory access}
\label{section:parallel-programming:shared-mem:protect}

In the tradition of Intel's original TBB version, I offer an abstraction over
critical sections aka exclusive code blocks.
That is blocks that can only be executed by one thread at a time.

To use it, you define a \texttt{BooleanSemaphore} object.
It should be global (or static) and is an abstraction around C++/TBB/OpenMP.
To lock it, you can do so manually, but actually I offer a \texttt{Lock} object
which locks the semaphore and automatically frees it once the object is
destroyed.

For details, please consult the documentation of the class \linebreak
\texttt{tarch::multicore::BooleanSemaphore}. The \texttt{multicore} namespace
also hosts variants to support recursive semaphores.

\begin{remark}
 It might be faster to use atomics if you want to protect the access to a single
 native variable. In this case, I recommend to stick to C++ atomics. The
 semaphore is a wrapper around C++ and OpenMP (so your code will work with
 both). If you use only C++ atomic, you should still be fine when you switch to
 OpenMP.
\end{remark}


\section{Distributed memory}

Distributed memory programming in Peano can be tricky, as we have a lot of
different threads per rank and the ranks run pretty asynchronously.
Therefore, it is always a bad idea to use blocking or synchronising routines. 
Here are a few recommendations:

\begin{enumerate}
  \item Do not use a blocking send or receive, even if routines are logically
  blocking.
  Instead, use a non-blocking routine and poll the termination via
  \texttt{MPI\_Test}.
  While you poll, invoke \texttt{receiveDanglingMessages()} on the services to
  allow \Peano\ to free the queue from other messages such as load balancing
  requests which might be in your way.
  \item Register your own routine as a service (see
  \texttt{tarch::services::Service}) and equip it with a
  \texttt{receiveDanglingMessages()} routine which probes the MPI queues. If a
  message of relevance is in the MPI queue, receive it. If you can't react
  directly, buffer it internally. This routine ensures that the MPI buffers
  can't overflow. Overflows force MPI to fall back from eager to rendeszouz
  communication which can lead to deadlocks.
  \item Do never impose any synchronisation unless absolutely required.
  Notably, try to avoid synchronisation within the action sets! In the ideal
  case, you implement a Jacobi-like data exchange: The action sets send stuff
  out.
  There's a service (see above) which receives the data. This data is then used
  by the action set (or another one) in the subsequent grid pass. 
  \item If you have a global operation, use a non-blocking collective and wait
  for the collective's test in the subsequent grid run-through.
\end{enumerate}

\begin{remark}
 If a message consists of multiple parts, you might have to fuse it into one
 message to avoid that the first part of a message is received and queue polling
 (on another thread) then grabs the second part misinterpreting it for the first
 part of yet another message.
\end{remark}


\subsection{Synchronisation}

I also offer a boolean semaphore accross all MPI ranks. This one is quite
expensive, as it has to lock both all threads and all ranks and thus send
quite some messages forth and back. 

The usage is similar to the shared memory lock, but you have to give the
semaphore a unique name (you might construct a string using the
\texttt{\_\_LINE\_\_} and \texttt{\_\_FILE\_\_} macro). 
It is the name that \Peano\ uses to ensure that the semaphore is unique accross
all ranks.
In this way, the semaphore is similar to an OpenMP named critical section.
I recommend to avoid MPI barrier whenever possible.
Details can be found in \texttt{tarch::mpi::Rank::barrier()}.



\subsection{Exchange global data between ranks (reduction)}
\label{section:parallel-programming:shared-mem:reductions}

If you reduce data betwen the ranks, I recommend that you use \Peano's wrappers.
They have two advantages: 
They automatically provide deadlock tracking, and they poll the MPI queues,
i.e.~if the code stalls as too many other messages for other purposes (global
load balancing, plotting, \ldots) drop in, and remove these messages from the
buffers such that MPI can make progress. The usage is detailed in
\texttt{tarch::mpi::Rank::allReduce}, but the code usually resembles:

\begin{code}
double myBuffer = _myVariable;

tarch::mpi::Rank::getInstance().allReduce(
  &myBuffer,
  &_myVariable,
  1, MPI_DOUBLE,
  MPI_MIN,
  [&]() -> void { tarch::services::ServiceRepository::getInstance().receiveDanglingMessages(); }
);
\end{code}


\section{Design patterns for global variables between action sets}

Realising global variables shared between action sets is in principle
impossible: 
Action sets are not persistent. 
They are created per traversal per tree.
Furthermore, you have absolutely no knowledge how many trees any rank hosts.
The only thing you know is that all trees will kind of start their traversal
around the same time.


However, there is a design pattern that you can use.
It relies on the observations that (i) all the \texttt{beginTraversal()} events
are triggered \emph{after} all the global trees have synchronised.
(ii) Tree 0 exists always.


\begin{itemize}
  \item There are two \emph{static} datasets per rank. This is the only valid
  data in-between any two grid traversals. One is the current data, one is the
  data from the previous sweep.
  \item Every action sets builds up its local current data and commits this 
  data to the static global dataset in \texttt{endTraversal()}. This local
  reduction has to be protected by a (rank-local) semaphore.
  \item In the main routine, each rank calls a routine which reduces the
  static current data set to rank 0. Rank 0 then broadcasts the reduced data and
  each rank rolls over the reduced data into the static old data set. The
  current data is cleared.
  \item Every action set copies the old data set into its local copy.
\end{itemize}


\begin{remark}
  These MPI/thread-global barriers can be very expensive and have the potential
  to lead to deadlocks/timeouts. Use them with care and localise their
  behaviour.
\end{remark}

